\chapter*{Abstract}
\addcontentsline{toc}{chapter}{Abstract}
\markboth{Abstract}{Abstract}
\begin{singlespace}

An image pair can now produce a dense 3D Gaussian reconstruction in one
network pass. A video instead creates overlapping predictions that must follow
loop-closure corrections. This thesis presents Splatt3R-SLAM, a monocular
system that turns local predictions into a persistent map. It gives each
local Gaussian map and its observation cameras the same keyframe anchor, so
pose corrections move both together. A separate process refines appearance
while MASt3R-SLAM estimates motion.

Head-only adaptation changes prediction, opacity attenuation changes
insertion, and photometric refinement uses sequence observations. Refinement
improves PSNR by $1.00$--$8.69$\,dB across four controls. Attenuation
improves LPIPS in ten of twelve tests, although confidence weighting offers
little advantage over uniform attenuation.

Under common non-keyframe reconstruction evaluation, the released head with
$300$\,s refinement reaches $23.02$\,dB mean PSNR on eight Replica scenes,
compared with $19.48$\,dB for Photo-SLAM. On TUM desk, its $13.95$\,dB is
close to MonoGS, which has lower LPIPS. Dense prediction provides a useful
initialisation but creates millions of primitives. Compact online mapping
remains the main open problem.
\end{singlespace}
